\section{Complete Human Adjudication and Skill Breakdown}
\label{app:human}

\subsection{Coverage and Review Procedure}

The review binds saved answer text and accepted scores from the current
305-task Ref experiment. A deterministic task-ID hash selects two tasks per
Skill, requiring at least eight usable model--condition cells out of ten.
This yields 18 source-scenario tasks and ten constructed tasks. The allocation
ensures Skill coverage rather than representing the benchmark's uneven
frequency distribution. The five model groups and both access conditions
match the main API comparison; GLM-5.3 retains its recorded historical
execution provenance.

All 280 sampled cells are accounted for: 249 substantive pairs, 28
reference-self controls, and three execution-failure controls. Every
substantive pair has a complete set of scores, textual evidence,
counterevidence where identified, and reasons. Controls retain their
predefined values and do not enter agreement statistics. No answers were
regenerated and no additional Judge requests were made for this analysis.

The initial scores and evidence were AI-proposed. Two expert reviewers
subsequently jointly reviewed and approved all 249 candidate label sets and
28 reference label sets, as confirmed by the study organizer. Affiliations
are anonymized. The resulting human-adjudicated labels are one shared set;
there are no separate independent reviewer scores. We therefore do not
estimate inter-rater reliability or claim an independent multi-rater design.

Review inputs preserve the task, applicable rubric, original anonymous A/B
order, and E-numbered delivery text, without displaying the saved numerical
Judge scores. Proposal preparation had prior exposure to aggregates and
some pilot cases; the joint human review is not documented as fully blinded.
Each reference is scored once and its labels are reused across candidates.
The data contain 1,345 candidate dimension scores and 152 unique reference
dimension scores (1,497 labels). The expanded comparison has 2,690 rows,
including reused references, not 2,690 independent human judgments.

\subsection{Comparisons and Denominators}

Both sides use the same applicable dimensions and continuous 0--10 scale.
For $D$ dimensions,
\[
\begin{aligned}
J_c&=\frac{1}{10D}\sum_d s_{c,d},\qquad
J_r=\frac{1}{10D}\sum_d s_{r,d},\\
R&=\frac{J_c}{J_c+J_r}.
\end{aligned}
\]
Ratios are calculated per pair before averaging. Parity at 0.5 is not
50\% accuracy; no pair has an undefined zero denominator. MAE and signed
bias compare paired values, with bias defined as adjudicated minus Judge.
Spearman correlation uses average ranks for ties and returns undefined for
constant inputs. Statistics are descriptive: repeated model outputs and
reference scores cluster within 28 tasks.

On 1,345 candidate dimensions, MAE is 0.992 on the 0--10 scale and
Spearman correlation is 0.270. Scores of at least nine occur on 1,163
dimensions under the Judge and 780 after adjudication; 234 adjudicated
dimensions still receive ten. Thus the review identifies localized
disagreements without imposing a ceiling on valid high scores.
Table~\ref{tab:ref-human-skill-metrics} reports every Skill rather than
selecting only favorable correlations.

\begin{table*}[t]
\centering\small
\setlength{\tabcolsep}{6pt}
\begin{tabular}{lrrrrrr}
\toprule
Skill & Pairs & $J_c$ MAE & $J_r$ MAE & $R$ MAE & $R$ $\rho$ & Direction matches\\
\midrule
\inputtable{tables/full_ref_skill_metrics.tex}
\bottomrule
\end{tabular}
\caption{Complete 14-Skill human--Judge breakdown, two task clusters per
row. MAEs are percentage points. Reference $J_r$ is expanded over pairs;
its reused human labels are not independent observations.}
\label{tab:ref-human-skill-metrics}
\end{table*}

\begin{table*}[t]
\centering\small
\begin{tabular}{lrrrrr}
\toprule
Task origin & Tasks & Pairs & Judge $J_c$ & Human $J_c$ & $R$ MAE\\
\midrule
\inputtable{tables/full_ref_source_metrics.tex}
\bottomrule
\end{tabular}
\caption{Task-origin breakdown. $J_c$ is on a 0--100 scale; $R$ MAE is
in percentage points. These coverage-sampled means do not replace the
305-task leaderboard.}
\label{tab:ref-human-source-metrics}
\end{table*}

\subsection{Reference Effects and Error Interpretation}

The same saved reference receives different automated scores across
candidate contexts: its within-task $J_r$ range averages 11.13 points and
reaches 25.00. Comparing each task's mean automated reference score against
its single adjudicated score gives a 6.88-point MAE over 28 unique references.
The absence of within-task variation in adjudicated references follows
from label reuse, not repeated-review stability.

Exact three-way preference agreement is 53.41\%. Treating scores within
one or two points of parity as ties changes agreement to 46.18\% or
62.65\%, respectively. These thresholds are sensitivity checks, not tuned
choices for the headline result. Component differences and weak pooled
ranking correlation explain why a small ratio MAE cannot alone validate
model ordering.

The evidence includes incorrect polynomial values and roots in a retrieval
task, inconsistent power calculations in an electromagnetic self-explanation
task, and a linear rather than square-root pendulum relation in an emergent
project. Other disagreements concern feasibility: mandatory emotion disclosure
despite stated privacy protections, or a delivery summary that does not
contain the claimed artifact. Reviewers assess the provided text rather
than unseen files and reduce related dimensions rather than erasing valid
teaching structure. References are subject to the same checks and are not
gold solutions. External factual verification is not exhaustive; this
study supports an auditable, Skill-specific diagnosis of scoring behavior,
not universal correctness guarantees.
